% Standard twirling, then partial twirls on successive outgoing qubits.
\begin{frame}{Iterative Pauli Twirling}
  \centering
  \begin{tikzpicture}[
      >={Latex[length=2.3mm]},
      process/.style={draw=CVQCInk,fill=white,thick,minimum width=1.5cm,
        minimum height=0.85cm,inner sep=3pt,font=\normalsize,text=CVQCInk},
      pad/.style={process,draw=CVQCTeal,fill=CVQCPale},
      attack/.style={process,draw=CVQCCoral,fill=CVQCCoralPale},
      wire/.style={thick,draw=CVQCInk},
      heading/.style={font=\normalsize\bfseries,text=CVQCInk},
      equation/.style={font=\large,text=CVQCInk}
    ]
    \path[use as bounding box] (-6.6,-3.15) rectangle (6.6,3.05);

    \only<1-2>{
      \node[heading] at (0,2.65) {Standard Pauli twirling};

      \node[pad] (enc) at (-2.6,1.55) {$\operatorname{OTP}$};
      \node[attack,minimum width=1.3cm] (attack) at (0,1.55) {$A$};
      \node[pad,minimum width=1.7cm] (dec) at (2.6,1.55)
        {$\operatorname{OTP}^{\dagger}$};
      \draw[wire] (-4.3,1.55) -- (enc.west);
      \draw[wire] (enc.east) -- (attack.west);
      \draw[wire] (attack.east) -- (dec.west);
      \draw[wire] (dec.east) -- (4.3,1.55);

      \only<2>{
        \draw[->,very thick,draw=CVQCTeal] (0,0.90) -- (0,-0.05)
          node[midway,right=8pt,font=\small,text=CVQCInk] {reduce};
        \node[attack,minimum width=1.3cm] (pauli) at (0,-0.80) {$Q$};
        \draw[wire] (-2.4,-0.80) -- (pauli.west);
        \draw[wire] (pauli.east) -- (2.4,-0.80);
        \node[equation] at (0,-2.15)
          {$A=\sum_Q\alpha_Q Q,\qquad \Pr[Q]=|\alpha_Q|^2$};
      }
    }

    \only<3-5>{
      \begin{scope}[yshift=0.30cm]

      \node[attack,minimum width=1.15cm] (secondattack) at (2.9,-0.65)
        {$A_2$};
      \node[pad,minimum width=1.7cm] (decsecond) at (5.0,-0.65)
        {$\operatorname{OTP}_2^{\dagger}$};

      \draw[wire] (secondattack.east) -- (decsecond.west);
      \draw[wire] (decsecond.east) -- (6.20,-0.65);

      \only<3,5>{
        \node[pad,minimum height=2.5cm] (encboth) at (-4.85,0.20)
          {$\operatorname{OTP}$};
        \node[attack,minimum width=1.6cm,minimum height=2.5cm,font=\small]
          (firstattack) at (-2.6,0.20)
          {$\alt<3>{A_1}{\mathrm{CNOT}_{2\to1}}$};
        \node[pad,minimum width=1.7cm] (decfirst) at (-0.15,1.05)
          {$\operatorname{OTP}_1^{\dagger}$};
        \foreach \y in {1.05,-0.65} {
          \draw[wire] (-6.15,\y) -- (encboth.west |- 0,\y);
          \draw[wire] (encboth.east |- 0,\y) -- (firstattack.west |- 0,\y);
        }
        \draw[wire] (firstattack.east |- 0,1.05) -- (decfirst.west);
        \draw[wire] (decfirst.east) -- (2.65,1.05);
        \draw[wire] (firstattack.east |- 0,-0.65) -- (secondattack.west);
      }
      \only<4>{
        \node[pad] (encremaining) at (-4.85,-0.65) {$\operatorname{OTP}_2$};
        \node[attack,minimum width=1.6cm] (firstpauli) at (-2.6,1.05) {$Q$};
        \node[attack,minimum width=1.6cm] (coefficient) at (-2.6,-0.65) {$B_Q$};
        \draw[wire] (-6.15,1.05) -- (firstpauli.west);
        \draw[wire] (firstpauli.east) -- (2.65,1.05);
        \draw[wire] (-6.15,-0.65) -- (encremaining.west);
        \draw[wire] (encremaining.east) -- (coefficient.west);
        \draw[wire] (coefficient.east) -- (secondattack.west);
      }
      \only<3-4>{
        \node[equation] at (0,-2.0) {$A_1=\sum_Q Q\otimes B_Q$};
      }
      \only<5>{
        \node[equation] at (0,-1.85)
          {$\mathrm{CNOT}_{2\to1}
            =I\otimes\ket{0}\!\bra{0}+X\otimes\ket{1}\!\bra{1}$};
      }
      \end{scope}
    }
  \end{tikzpicture}

  \note<1>{
    \footnotesize
    Start with the standard Pauli-twirling identity. In this first toy
    picture, OTP means a uniform secret Pauli pad, X to the u times Z to
    the v, and A is a key-independent unitary attack. This differs from
    the X-to-the-x T-to-the-t notation for the MBQC input pads: the actual
    MBQC analysis first uses the uniformity-preserving change of keys to
    expose the consumed Pauli pad bits. Expand A in the Pauli basis,
    with scalar coefficients alpha indexed by the Paulis Q. The wire may
    represent a register, and the identity also preserves an external
    reference. The decryption is part of an analysis experiment; the
    picture is not assigning extra quantum operations to our verifier.
  }
  \note<2>{
    \footnotesize
    Average over the hidden Pauli key. Conjugation multiplies each Pauli
    term by a key-dependent sign; the cross terms between distinct Paulis
    vanish in the density operator. The decrypted channel is therefore
    the classical mixture that applies Q with probability absolute alpha-Q
    squared. This is an equality after averaging the key, not equality
    between A and a single Pauli as operators or for each fixed key.
    The scalar probabilities here describe the unitary toy example. With
    arbitrary private memory, expansion coefficients can themselves be
    operators on that memory. This leads to the partial twirl next.
    Source: the manuscript's Technical Overview, Pauli-twirling discussion.
  }
  \note<3>{
    \footnotesize
    Now peel off the pad on one outgoing qubit per round, keeping the
    other qubit encrypted. The initial box pads both qubits independently.
    A1 acts jointly; decrypt and release the first outgoing qubit, then allow A2 to act on
    the remaining ciphertext before its later decryption. The picture is
    a two-qubit toy; additional qubits, private memory, and transcript
    dependence are implicit. Its later measurement in MBQC is suppressed
    here. Expand only on the first qubit as
    $A_1=\sum_Q Q\otimes B_Q$, where Q ranges over I,X,Y,Z. The $B_Q$
    coefficients are operators on the remaining system. Twirling removes
    interference between different Q labels, while their coefficient
    operators and branch weights must be carried through later rounds.
    In MBQC, this averaging is justified by the consumed hidden pad
    coordinates. The full proof also retains the prover's memory and live
    verifier keys. These labels are analysis indices, not public outcomes.
  }
  \note<4>{
    \footnotesize
    Average over the consumed first-qubit Pauli key. Its encryption and
    decryption disappear, leaving Q on the first wire and the corresponding
    coefficient $B_Q$ on the second. The second qubit remains encrypted;
    its pad is consumed in the later round. The figure displays one term
    of the resulting Kraus sum. Sum over Q to obtain the complete map;
    $B_Q$ carries the branch weight, so this is not a uniform choice among
    normalized operators. Keep the state subnormalized through A2 and the
    remaining decryption. Q is an analysis index; it is not revealed to the
    prover or the verifier. Next, instantiate A1 with CNOT to see why its
    coefficient operators need not be unitary.
  }
  \note<5>{
    \footnotesize
    The paper's counterexample is CNOT from qubit 2 to qubit 1. Its partial
    expansion on the first qubit has I and X there, but the coefficients
    on the second qubit are the projectors onto zero and one. Those
    projectors can shrink a vector's norm. After a partial twirl, we must
    therefore carry subnormalized branch states through A2 and the later
    twirl, rather than normalizing each branch and losing its weight.
    A later expansion can decompose the remaining projector, for example
    $\ket{0}\!\bra{0}=(I+Z)/2$. Even a globally unitary attack thus
    generates nonunitary coefficient operators during the iteration.
    Track the sum of branch traces throughout the iteration.
    Source: the Technical Overview and Real-Copy Protocol,
    partial-twirl lemma and iterative analysis. Next, apply this mechanism
    to the actual MBQC measurement gadget.
  }
\end{frame}
